\documentclass[../../../include/open-logic-section]{subfiles}

\begin{document}

\olfileid{his}{set}{pathology}
\olsection{విచిత్ర నిర్మాణాలు}

అయితే \emph{పరిమితి} నిర్వచనం కొన్ని ``విచిత్రమైన'' నిర్మాణాలకూ అవకాశం ఇస్తుందని తేలింది.

దాదాపు 1830లలో బోల్జానో, \emph{ప్రతిచోటా అవిచ్ఛిన్నం} అయినా \emph{ఎక్కడా అవకలనీయంకాని} ఒక ప్రమేయాన్ని కనుగొన్నారు. (దురదృష్టవశాత్తూ బోల్జానో దాన్ని ప్రచురించలేదు; అదే భావనను వైయర్‌స్ట్రాస్ స్వతంత్రంగా కనుగొనడంతో 1872లో గణితశాస్త్రజ్ఞులు దానిని తొలిసారిగా చూశారు.)\footnote{ఈ చరిత్ర \href{http://en.wikipedia.org/wiki/Weierstrass_function}{వైయర్‌స్ట్రాస్ ప్రమేయం} గురించి వికీపీడియా వ్యాసంలోని అత్యంత విస్తృత పాదసూచికల్లో నమోదు అయింది.} ఇది ఎంత చెప్పినా ఆశ్చర్యకరమే. $|x|$ వంటి, ప్రతిచోటా అవిచ్ఛిన్నంగా ఉండి ఒక నిర్దిష్ట బిందువు వద్ద అవకలనీయంకాని ప్రమేయాలు కనుగొనడం సులభం. కానీ ప్రతిచోటా అవిచ్ఛిన్నంగా ఉండి \emph{ఎక్కడా} అవకలనీయంకాని ప్రమేయం పూర్తిగా భిన్నమైనది. అటువంటి ప్రమేయాన్ని ఎలా గీయగలరో ఒక్కసారి ఆలోచించండి. అది అవిచ్ఛిన్నంగా ఉండాలంటే కలాన్ని కాగితం మీద నుంచి ఎప్పుడూ తీయకుండా గీయగలగాలి; కానీ ఎక్కడా అవకలనీయంకాకూడదంటే కలం దిశను నిరంతరం అకస్మాత్తుగా మార్చాల్సి ఉంటుంది.

ఇంకా కొన్ని ``విచిత్ర నిర్మాణాలు'' వచ్చాయి. 1874 జనవరి 5న కాంటర్ డెడెకిండ్‌కు ఒక లేఖలో ఈ సమస్యను ప్రతిపాదించారు:
\begin{quote}
ఒక ఉపరితలానికి (దాని అంచులతో సహా ఒక చతురస్రం అనుకుందాం), ఒక రేఖకు (దాని చివరి బిందువులతో సహా ఒక సరళరేఖ అనుకుందాం) మధ్య ఒకటి-ఒకటి అనుగుణ్యత ఉండగలదా? అంటే ఉపరితలంలోని ప్రతి బిందువుకు రేఖలో ఒక బిందువు, తిరిగి రేఖలోని ప్రతి బిందువుకు ఉపరితలంలో ఒక బిందువు అనుగుణంగా ఉండగలదా?

ఈ ప్రశ్నకు సమాధానం ఇప్పటికీ నాకు చాలా కష్టమైనదిగా అనిపిస్తోంది---ఇక్కడ కూడా \emph{కాదు} అని అనాలని బలంగా అనిపిస్తున్నందువల్ల, నిరూపణ దాదాపు అనవసరమని భావించాలని ఉంది. [\citealt{Gouvea2011}లో ఉటంకింపు]
\end{quote}
కానీ 1877లో తాను తప్పుగా భావించానని కాంటర్ నిరూపించారు. వాస్తవానికి రేఖలోనూ చతురస్రంలోనూ బిందువుల సంఖ్య సరిగ్గా సమానమే. 1877 జూన్ 29న డెడెకిండ్‌కు ``\emph{je le vois, mais je ne le crois pas}'' అని రాశారు; అంటే ``నాకు అది కనిపిస్తోంది, కానీ నమ్మలేకపోతున్నాను''. గణితశాస్త్రపు ``సాధారణంగా అంగీకరించిన చరిత్ర''లో దీన్ని, ఆ కాలపు గణితశాస్త్రజ్ఞులకు ఈ కొత్త ఫలితాలు ఎంత \emph{నిజంగానే నమ్మశక్యం కానివో} తెలిపేదిగా తరచుగా తీసుకున్నారు. (ఆ ఉత్తరప్రత్యుత్తరాలు \citet{Gouvea2011}లో ఉన్నాయి; \olref[his][set][mythology]{sec}లో మళ్లీ వాటిని చూస్తాం. కాంటర్ నిరూపణ రూపురేఖలు \olref[his][set][cantorplane]{sec}లో ఉన్నాయి.)

కాంటర్ ఫలితంతో ప్రేరేపితుడైన పియానో, ఒక రేఖఖండాన్ని తలంలోని చతురస్రంపై \emph{అవిచ్ఛిన్నంగా} పంపడం సాధ్యమా అని పరిశీలించడం మొదలుపెట్టారు. అది \emph{స్థలాన్ని నింపే వక్రరేఖ} అవుతుంది. \citeyear{Peano1890}లో పియానో అలాంటి వక్రరేఖను నిర్మించారు.
\sourcecorrection{OLTEHISSETPATH-001}{మూలం ఇక్కడ ``సున్నితంగా'' తలంపై పంపడం అని చెప్పి, పియానో వక్రరేఖనే అటువంటి ఉదాహరణగా ఇచ్చింది. సున్నితమైన (నిరంతర వ్యుత్పన్నం గల) వక్రరేఖ చతురస్రాన్ని నింపలేదు; పియానో నిర్మాణం రేఖఖండం నుంచి చతురస్రానికి అవిచ్ఛిన్న పటం. ఆ భేదాన్ని స్పష్టం చేశాం.}
ఇది నిజంగా సహజ అంచనాలకు విరుద్ధం: యూక్లిడ్ రేఖను ``వెడల్పు లేని పొడవు''గా నిర్వచించారు (పుస్తకం I, నిర్వచనం 2); కానీ రేఖను తగిన విధంగా మడతపెట్టడం ద్వారా దాని పొడవును వెడల్పుగా మార్చవచ్చని పియానో చూపారు. \citeyear{Hilbert1891}లో హిల్బర్ట్ కొంచెం సులభంగా ఊహించగల స్థలాన్ని నింపే వక్రరేఖను, దాన్ని వివరించే కొన్ని చిత్రాలతో సహా, వర్ణించారు. ఆ వక్రరేఖను దశలవారీగా నిర్మిస్తారు; నిర్మాణంలోని తొలి ఆరు దశలు ఇవి:
\begin{center}
\begin{tikzpicture}
\begin{scope}[xshift=20pt, yshift=20pt]
	\draw [oldiagcolorC, l-system={Hilbert curve, step=40pt, angle=90, axiom=L, order=1}] lindenmayer system;
\end{scope}
\begin{scope}[xshift=110pt, yshift=10pt]
	\draw [oldiagcolorC, l-system={Hilbert curve, step=20pt, angle=90, axiom=L, order=2}] lindenmayer system;
\end{scope}
\begin{scope}[xshift=205pt, yshift=5pt]
	\draw [oldiagcolorC, l-system={Hilbert curve, step=10pt, angle=90, axiom=L, order=3}] lindenmayer system;
\end{scope}
\begin{scope}[xshift=2.5pt, yshift=-97.5pt]
	\draw [oldiagcolorC, l-system={Hilbert curve, step=5pt, angle=90, axiom=L, order=4}] lindenmayer system;
\end{scope}
\begin{scope}[xshift=101.25pt, yshift=-98.75pt]
	\draw [oldiagcolorC, l-system={Hilbert curve, step=2.5pt, angle=90, axiom=L, order=5}] lindenmayer system;
\end{scope}
\begin{scope}[xshift=200.625pt, yshift=-99.375pt]
	\draw [oldiagcolorC, l-system={Hilbert curve, step=1.25pt, angle=90, axiom=L, order=6}] lindenmayer system;
\end{scope}
\draw[gray] (0pt,0pt) rectangle (80pt, 80pt);
\draw[gray] (100pt,0pt) rectangle (180pt, 80pt);
\draw[gray] (200pt,0pt) rectangle (280pt, 80pt);
\draw[gray] (0pt,-100pt) rectangle (80pt, -20pt);
\draw[gray] (100pt,-100pt) rectangle (180pt, -20pt);
\draw[gray] (200pt,-100pt) rectangle (280pt, -20pt);
\end{tikzpicture}  
\end{center}
పరిమితి దశలో---అప్పటికే కచ్చిత నిర్వచనం పొందిన భావన ఇది---మొత్తం చతురస్రం దట్టమైన \tetoken{ఎరుపు}{!!{colorC}} రంగుతో నిండిపోతుంది. సందర్భవశాత్తూ, హిల్బర్ట్ వక్రరేఖ ప్రతిచోటా అవిచ్ఛిన్నం, కానీ ఎక్కడా అవకలనీయం కాదు; ఎందుకంటే అనంత పరిమితి దశలో ఆ ప్రమేయం ప్రతి క్షణం అకస్మాత్తుగా దిశ మారుస్తుందని సహజంగా భావించవచ్చు. (హిల్బర్ట్ నిర్మాణాన్ని \olref[his][set][hilbertcurve]{sec}లో మరింత వివరంగా చూస్తాం.)

మంచికో చెడుకో, ఈ ``విచిత్ర'' జ్యామితీయ నిర్మాణాలను జ్యామితీయ అంతఃప్రజ్ఞపై ఆధారపడటాన్ని సందేహించడానికి కారణంగా తీసుకున్నారు. అవి సమితి సిద్ధాంతంలో పరాకాష్ఠకు చేరే కొత్త గణిత పద్ధతికి దాదాపు \emph{ప్రచార సాధనాలు} అయ్యాయి. ఆ రంగపు తరువాతి పురాణ నిర్మాణంలో, ఈ ఫలితాలు పూర్తిగా కచ్చితమైనవీ, పూర్తిగా దిగ్భ్రాంతికరమైనవీ అని తరచుగా పునరుక్తి చేశారు. అందువల్ల అవి రెండు పనులు చేశాయి: జ్యామితీయ అంతఃప్రజ్ఞపై ఆధారపడవద్దనే హెచ్చరికగా, కొత్త ఆలోచనా విధానాల ఫలప్రదతకు నిదర్శనంగా.

\end{document}
